% Section 2.4, from lectures/L02/appendix/c_exercises.md.

\section{Exercises}
\label{c2:sec:exercises}\label{c2:app:c}

\begin{aside}
\textbf{How to work these.} These exercises reinforce the concepts from \secref{c2:sec:classes1} and
\secref{c2:sec:classes2}. Build each program with a Makefile like the one in
\secref{c1:sec:makefile}; for a class split across header and source files, use the
\file{include}/\file{source} layout and Makefile from \secref{c1:sec:tips}. Check your output against
the example output each exercise gives.

The solutions are published in the companion repository, one program per set under
\file{lectures/L02/appendix/solutions}, each with its Makefile. Try each exercise yourself before
looking at them.
\end{aside}

% ------------------------------------------------------------------------------------------
\exerciseset{First Class}
\label{c2:set:1}

\exercise{Simple LED class}{Code}
\label{c2:ex:1.1}
In a new program, create a class \code{Led} in the namespace \code{driver::gpio}:

\begin{cppcode}
namespace driver::gpio
{
class Led final
{
public:

private:
};
} // namespace driver::gpio
\end{cppcode}

\task{Tasks}
\begin{parts}
  \item Add two private member variables:
  \begin{itemize}
    \item The first member variable shall:
    \begin{itemize}
      \item Be named \code{myPin}.
      \item Represent the pin number to which the LED is connected.
      \item Have the type \code{std::uint8_t}.
      \item Be readable only after initialization.
    \end{itemize}
    \item The second member variable shall:
    \begin{itemize}
      \item Be named \code{myState}.
      \item Represent the state of the LED (\code{true/false}).
      \item Be \code{true} when the LED is enabled.
    \end{itemize}
  \end{itemize}
  \item Add a constructor that:
  \begin{itemize}
    \item Takes a pin number (\code{std::uint8_t}).
    \item Takes an initial state (\code{bool}) with the default value \code{false}.
    \item Uses an initialization list to initialize the member variables.
    \item Is marked \code{explicit} and \code{noexcept}.
  \end{itemize}
  \item Add a destructor and mark it \code{noexcept} and \code{default}.
  \item Implement the following methods directly in the class:
  \begin{itemize}
    \item The method \code{isOn()} shall:
    \begin{itemize}
      \item Return \code{true} if the LED is enabled, otherwise \code{false}.
      \item Be marked \code{const}, \code{noexcept}, and \code{[[nodiscard]]}.
    \end{itemize}
    \item The method \code{on()} shall:
    \begin{itemize}
      \item Take no parameters.
      \item Set the LED state to \code{true}.
      \item Return no value.
      \item Be marked \code{noexcept}.
    \end{itemize}
    \item The method \code{off()} shall:
    \begin{itemize}
      \item Take no parameters.
      \item Set the LED state to \code{false}.
      \item Return no value.
      \item Be marked \code{noexcept}.
    \end{itemize}
    \item The method \code{toggle()} shall:
    \begin{itemize}
      \item Toggle the LED state.
      \item Take no parameters.
      \item Return no value.
      \item Be marked \code{noexcept}.
    \end{itemize}
  \end{itemize}
  \item Delete the following methods:
  \begin{itemize}
    \item The default constructor.
    \item The copy constructor.
    \item The move constructor.
    \item The copy assignment operator.
    \item The move assignment operator.
  \end{itemize}
  \item Create an instance of the class:
\begin{cppcode}
constexpr std::uint8_t ledPin{13U};
driver::gpio::Led led{ledPin};
\end{cppcode}
  Test the class by:
  \begin{itemize}
    \item Printing the state.
    \item Calling \code{on()}.
    \item Printing again.
    \item Calling \code{toggle()}.
    \item Printing once more.
  \end{itemize}
  Use \code{std::printf()} from \code{<cstdio>} for the output.
\end{parts}
Example output:

\begin{console}
Initial state: Off
After on(): On
After toggle(): Off
\end{console}

\exercise{Button class}{Code}
\label{c2:ex:1.2}
In the same program, create a class \code{Button} in the namespace \code{driver::gpio}:

\begin{cppcode}
namespace driver::gpio
{
class Button final
{
public:

private:
};
} // namespace driver::gpio
\end{cppcode}

\task{Tasks}
\begin{parts}
  \item Add two private member variables:
  \begin{itemize}
    \item The first member variable shall:
    \begin{itemize}
      \item Be named \code{myPin}.
      \item Represent the pin number to which the button is connected.
      \item Have the type \code{std::uint8_t}.
      \item Be readable only after initialization.
    \end{itemize}
    \item The second member variable shall:
    \begin{itemize}
      \item Be named \code{myPressed}.
      \item Represent the state of the button (\code{true/false}).
      \item Be \code{true} when the button is pressed.
    \end{itemize}
  \end{itemize}
  \item Add a constructor that:
  \begin{itemize}
    \item Takes a pin number (\code{std::uint8_t}).
    \item Uses an initialization list to initialize the member variables.
    \item Is marked \code{explicit} and \code{noexcept}.
  \end{itemize}
  \item Add a destructor and mark it \code{noexcept} and \code{default}.
  \item Implement the following methods directly in the class:
  \begin{itemize}
    \item The method \code{pin()} shall:
    \begin{itemize}
      \item Return the pin number.
      \item Be marked \code{const}, \code{noexcept}, and \code{[[nodiscard]]}.
    \end{itemize}
    \item The method \code{isPressed()} shall:
    \begin{itemize}
      \item Return \code{true} if the button is pressed, otherwise \code{false}.
      \item Be marked \code{const}, \code{noexcept}, and \code{[[nodiscard]]}.
    \end{itemize}
    \item The method \code{setPressed()} shall:
    \begin{itemize}
      \item Take an argument of type \code{bool}.
      \item Update the internal state of the button.
      \item Return no value.
      \item Be marked \code{noexcept}.
    \end{itemize}
  \end{itemize}
  \item Delete the following methods:
  \begin{itemize}
    \item The default constructor.
    \item The copy constructor.
    \item The move constructor.
    \item The copy assignment operator.
    \item The move assignment operator.
  \end{itemize}
  \item Create an instance:
\begin{cppcode}
constexpr std::uint8_t buttonPin{2U};
driver::gpio::Button button{buttonPin};
\end{cppcode}
  Simulate a button press:
\begin{cppcode}
button.setPressed(true);
\end{cppcode}
  Print the pin number and whether the button is pressed using \code{std::printf()} from
  \code{<cstdio>}.
\end{parts}
Example output:

\begin{console}
Button pin: 2
Pressed: Yes
\end{console}

% ------------------------------------------------------------------------------------------
\exerciseset{Interaction Between Classes}
\label{c2:set:2}

\exercise{Control LED with a button}{Code}
\label{c2:ex:2.1}
Use the classes from the previous exercises.

\task{Tasks}
\begin{parts}
  \item Create an LED on pin \code{13}.
  \item Create a button on pin \code{2}.
  \item If the button is not pressed, the LED should be off.
  \item If the button is pressed, the LED should turn on.
  \item Print the LED state.
\end{parts}
Example output:

\begin{console}
Button released -> LED: Off
Button pressed  -> LED: On
\end{console}

\exercise{Blink function}{Code}
\label{c2:ex:2.2}
Create a loop that toggles the LED state.

\task{Tasks}
\begin{parts}
  \item Turn the LED off, then run a loop for six iterations.
  \item In each iteration:
  \begin{itemize}
    \item Call \code{toggle()}.
    \item Print the LED state.
  \end{itemize}
\end{parts}
Example output:

\begin{console}
Blink 1: On
Blink 2: Off
Blink 3: On
Blink 4: Off
Blink 5: On
Blink 6: Off
\end{console}

% ------------------------------------------------------------------------------------------
\exerciseset{Class Split Across Multiple Files}
\label{c2:set:3}

\exercise{Buzzer class}{Code}
\label{c2:ex:3.1}
In this exercise, you will create a class \code{Buzzer} in the namespace \code{driver}. The class
should be split across multiple files. Therefore, create the following directory structure:

\begin{console}
Makefile
include/
    driver/
        buzzer.hpp
source/
    driver/
        buzzer.cpp
    main.cpp
\end{console}

\note
\begin{itemize}
  \item All method declarations shall be written in \file{driver/buzzer.hpp}.
  \item All method definitions shall be implemented in \file{driver/buzzer.cpp}.
\end{itemize}

In the header file \file{driver/buzzer.hpp}, the class shall be declared according to the structure
below:

\begin{cppcode}
#pragma once

#include <cstdint>

namespace driver
{
class Buzzer final
{
public:

private:
};
} // namespace driver
\end{cppcode}

\task{Tasks}
\begin{parts}
  \item Add two private member variables:
  \begin{itemize}
    \item The first member variable shall:
    \begin{itemize}
      \item Represent the pin number to which the buzzer is connected.
      \item Have the type \code{std::uint8_t}.
      \item Be named \code{myPin}.
      \item Be readable only after initialization.
    \end{itemize}
    \item The second member variable shall:
    \begin{itemize}
      \item Represent the state of the buzzer (\code{true}/\code{false}).
      \item Be named \code{myEnabled}.
      \item Be \code{true} when the buzzer is enabled.
    \end{itemize}
  \end{itemize}
  \item Add a constructor in the header file that:
  \begin{itemize}
    \item Takes a pin number (\code{std::uint8_t}).
    \item Takes an initial state (\code{bool}) with the default value \code{false}.
    \item Is marked \code{explicit} and \code{noexcept}.
  \end{itemize}
  \item Add a destructor in the header file:
  \begin{itemize}
    \item Mark it \code{noexcept}.
    \item Implement it in \file{driver/buzzer.cpp} so that the output below is generated when the
      buzzer is destroyed:
  \end{itemize}
\begin{console}
Releasing resources allocated for buzzer at pin 8!
\end{console}
  \note In the example output above, it is assumed that the buzzer is connected to pin 8.
  \item Declare the following methods in the header file and implement them in
    \file{driver/buzzer.cpp}:
  \begin{itemize}
    \item The method \code{pin()} shall:
    \begin{itemize}
      \item Return the pin number to which the buzzer is connected.
      \item Be marked \code{const}, \code{noexcept}, and \code{[[nodiscard]]}.
    \end{itemize}
    \item The method \code{isEnabled()} shall:
    \begin{itemize}
      \item Return \code{true} if the buzzer is enabled, otherwise \code{false}.
      \item Take no parameters.
      \item Be marked \code{const}, \code{noexcept}, and \code{[[nodiscard]]}.
    \end{itemize}
    \item The method \code{enable()} shall:
    \begin{itemize}
      \item Set the buzzer state to \code{true}.
      \item Take no parameters.
      \item Return no value.
      \item Be marked \code{noexcept}.
    \end{itemize}
    \item The method \code{disable()} shall:
    \begin{itemize}
      \item Set the buzzer state to \code{false}.
      \item Take no parameters.
      \item Return no value.
      \item Be marked \code{noexcept}.
    \end{itemize}
    \item The method \code{toggle()} shall:
    \begin{itemize}
      \item Toggle the buzzer state.
      \item Take no parameters.
      \item Return no value.
      \item Be marked \code{noexcept}.
    \end{itemize}
  \end{itemize}
  \item Delete the following methods in the header file:
  \begin{itemize}
    \item The default constructor.
    \item The copy constructor.
    \item The move constructor.
    \item The copy assignment operator.
    \item The move assignment operator.
  \end{itemize}
  \item In \file{main.cpp}, create an instance of the class:
\begin{cppcode}
constexpr std::uint8_t buzzerPin{8U};
driver::Buzzer buzzer{buzzerPin};
\end{cppcode}
  Test the class by:
  \begin{itemize}
    \item Printing the state.
    \item Calling \code{enable()}.
    \item Printing again.
    \item Calling \code{toggle()}.
    \item Printing again.
  \end{itemize}
  Use \code{std::printf()} from \code{<cstdio>} for the output.
\end{parts}
Example output:

\begin{console}
Initial state: Disabled
After enable(): Enabled
After toggle(): Disabled
Releasing resources allocated for buzzer at pin 8!
\end{console}

% ------------------------------------------------------------------------------------------
\exerciseset{Timer Class}
\label{c2:set:4}

\exercise{Timer}{Code}
\label{c2:ex:4.1}
In this exercise, you will create a class \code{Timer} in the namespace \code{driver}.

The class should be split across multiple files. Therefore, create the following files in the same
project as the previous exercise:

\begin{console}
include/driver/timer.hpp
source/driver/timer.cpp
\end{console}

\note
\begin{itemize}
  \item All method declarations shall be written in \file{driver/timer.hpp}.
  \item All method definitions shall be implemented in \file{driver/timer.cpp}.
\end{itemize}

In the header file, the class shall be declared according to the structure below:

\begin{cppcode}
#pragma once

#include <cstdint>

namespace driver
{
class Timer final
{
public:

private:
};
} // namespace driver
\end{cppcode}

\task{Tasks}
\begin{parts}
  \item Add four private member variables:
  \begin{itemize}
    \item The first member variable shall:
    \begin{itemize}
      \item Be named \code{myTimeout_ms}.
      \item Represent the timeout duration in milliseconds.
      \item Have the type \code{std::uint32_t}.
      \item Be readable only after initialization.
    \end{itemize}
    \item The second member variable shall:
    \begin{itemize}
      \item Be named \code{myCounter_ms}.
      \item Represent the internal timer counter.
      \item Have the type \code{std::uint32_t}.
    \end{itemize}
    \item The third member variable shall:
    \begin{itemize}
      \item Be named \code{myRunning}.
      \item Represent whether the timer is active or not (\code{true/false}).
      \item Be \code{true} when the timer is running.
    \end{itemize}
    \item The fourth member variable shall:
    \begin{itemize}
      \item Be named \code{myInitialized}.
      \item Represent whether the timer is initialized or not (\code{true/false}).
      \item Be \code{true} when the timer is initialized.
    \end{itemize}
  \end{itemize}
  \item Add a constructor in the header file that:
  \begin{itemize}
    \item Takes a timeout value in milliseconds (\code{std::uint32_t}).
    \item Uses an initialization list to initialize the member variables.
    \item Initializes the counter to \code{0}.
    \item Initializes the timer as stopped.
    \item Marks the timer as initialized if the timeout value is valid (\code{> 0}).
    \item Is marked \code{explicit} and \code{noexcept}.
  \end{itemize}
  \item Add a destructor in the header file:
  \begin{itemize}
    \item Mark it \code{noexcept}.
    \item Implement it in \file{driver/timer.cpp} so that:
    \begin{itemize}
      \item The timer is stopped if it was running.
      \item If the timer was running and is stopped, the output below should be printed:
    \end{itemize}
  \end{itemize}
\begin{console}
Stopping timer before deletion!
\end{console}
  \item Declare the following methods in the header file and implement them in
    \file{driver/timer.cpp}:
  \begin{itemize}
    \item The method \code{timeout_ms()} shall:
    \begin{itemize}
      \item Return the configured timeout in milliseconds.
      \item Take no parameters.
      \item Be marked \code{const}, \code{noexcept}, and \code{[[nodiscard]]}.
    \end{itemize}
    \item The method \code{isRunning()} shall:
    \begin{itemize}
      \item Return \code{true} if the timer is running, otherwise \code{false}.
      \item Take no parameters.
      \item Be marked \code{const}, \code{noexcept}, and \code{[[nodiscard]]}.
    \end{itemize}
    \item The method \code{isInitialized()} shall:
    \begin{itemize}
      \item Return \code{true} if the timer is initialized, otherwise \code{false}.
      \item Take no parameters.
      \item Be marked \code{const}, \code{noexcept}, and \code{[[nodiscard]]}.
    \end{itemize}
    \item The method \code{start()} shall:
    \begin{itemize}
      \item Start the timer if the timer is initialized.
      \item Take no parameters.
      \item Return no value.
      \item Be marked \code{noexcept}.
    \end{itemize}
    \item The method \code{stop()} shall:
    \begin{itemize}
      \item Stop the timer if the timer is initialized.
      \item Take no parameters.
      \item Return no value.
      \item Be marked \code{noexcept}.
    \end{itemize}
    \item The method \code{toggle()} shall:
    \begin{itemize}
      \item Toggle the timer if the timer is initialized.
      \item Take no parameters.
      \item Return no value.
      \item Be marked \code{noexcept}.
    \end{itemize}
    \item The method \code{tick()} shall:
    \begin{itemize}
      \item Simulate that 1 millisecond has passed.
      \item Increment the counter if the timer is running; otherwise do nothing.
      \item Return no value.
      \item Be marked \code{noexcept}.
    \end{itemize}
    \item The method \code{hasTimedOut()} shall:
    \begin{itemize}
      \item Return \code{true} when the counter has reached the timeout value, otherwise
        \code{false}.
      \item Reset the counter to \code{0} when a timeout occurs.
      \item Be marked \code{noexcept} and \code{[[nodiscard]]}.
    \end{itemize}
  \end{itemize}
  \item Delete the following methods:
  \begin{itemize}
    \item The default constructor.
    \item The copy constructor.
    \item The move constructor.
    \item The copy assignment operator.
    \item The move assignment operator.
  \end{itemize}
  \item In \file{main.cpp}, create an instance of the class:
\begin{cppcode}
constexpr std::uint32_t timeout_ms{1000U};
driver::Timer timer{timeout_ms};
\end{cppcode}
  Test the class by:
  \begin{itemize}
    \item Starting the timer.
    \item Running a loop for 3000 iterations. In each iteration:
    \begin{itemize}
      \item Call \code{tick()}.
      \item Check \code{hasTimedOut()}.
      \item Print a message if a timeout has occurred:
    \end{itemize}
  \end{itemize}
\begin{console}
Timeout after 1000 ms!
\end{console}
  \note In the example output above, it is assumed that the timeout is set to 1000 milliseconds.

  Use \code{std::printf()} from \code{<cstdio>} for the output above.
\end{parts}
Example output:

\begin{console}
Timeout after 1000 ms!
Timeout after 1000 ms!
Timeout after 1000 ms!
Stopping timer before deletion!
\end{console}
